\documentclass{amsart}
% For dealing with references we use the comment environment
\usepackage{verbatim}
\newenvironment{reference}{\comment}{\endcomment}
%\newenvironment{reference}{}{}
\newenvironment{slogan}{\comment}{\endcomment}
\newenvironment{history}{\comment}{\endcomment}

% For commutative diagrams we use Xy-pic
\usepackage[all]{xy}

% We use 2cell for 2-commutative diagrams.
\xyoption{2cell}
\UseAllTwocells

% We use multicol for the list of chapters between chapters
\usepackage{multicol}

% This is generally recommended for better output
\usepackage{lmodern}
\usepackage[T1]{fontenc}

% For cross-file-references

% Package for hypertext links:
\usepackage{hyperref}

% For any local file, say "hello.tex" you want to link to please

% Theorem environments.
%
\theoremstyle{plain}
\newtheorem{theorem}[subsection]{Theorem}
\newtheorem{proposition}[subsection]{Proposition}
\newtheorem{lemma}[subsection]{Lemma}

\theoremstyle{definition}
\newtheorem{definition}[subsection]{Definition}
\newtheorem{example}[subsection]{Example}
\newtheorem{exercise}[subsection]{Exercise}
\newtheorem{situation}[subsection]{Situation}

\theoremstyle{remark}
\newtheorem{remark}[subsection]{Remark}
\newtheorem{remarks}[subsection]{Remarks}

\numberwithin{equation}{subsection}

% Macros
%
\def\lim{\mathop{\mathrm{lim}}\nolimits}
\def\colim{\mathop{\mathrm{colim}}\nolimits}
\def\Spec{\mathop{\mathrm{Spec}}}
\def\Hom{\mathop{\mathrm{Hom}}\nolimits}
\def\Ext{\mathop{\mathrm{Ext}}\nolimits}
\def\SheafHom{\mathop{\mathcal{H}\!\mathit{om}}\nolimits}
\def\SheafExt{\mathop{\mathcal{E}\!\mathit{xt}}\nolimits}
\def\Sch{\mathit{Sch}}
\def\Mor{\mathop{\mathrm{Mor}}\nolimits}
\def\Ob{\mathop{\mathrm{Ob}}\nolimits}
\def\Sh{\mathop{\mathit{Sh}}\nolimits}
\def\NL{\mathop{N\!L}\nolimits}
\def\CH{\mathop{\mathrm{CH}}\nolimits}
\def\proetale{{pro\text{-}\acute{e}tale}}
\def\etale{{\acute{e}tale}}
\def\QCoh{\mathit{QCoh}}
\def\Ker{\mathop{\mathrm{Ker}}}
\def\Im{\mathop{\mathrm{Im}}}
\def\Coker{\mathop{\mathrm{Coker}}}
\def\Coim{\mathop{\mathrm{Coim}}}

% Boxtimes
%
\DeclareMathSymbol{\boxtimes}{\mathbin}{AMSa}{"02}

%
% Macros for moduli stacks/spaces
%
\def\QCohstack{\mathcal{QC}\!\mathit{oh}}
\def\Cohstack{\mathcal{C}\!\mathit{oh}}
\def\Spacesstack{\mathcal{S}\!\mathit{paces}}
\def\Quotfunctor{\mathrm{Quot}}
\def\Hilbfunctor{\mathrm{Hilb}}
\def\Curvesstack{\mathcal{C}\!\mathit{urves}}
\def\Polarizedstack{\mathcal{P}\!\mathit{olarized}}
\def\Complexesstack{\mathcal{C}\!\mathit{omplexes}}
% \Pic is the operator that assigns to X its picard group, usage \Pic(X)
% \Picardstack_{X/B} denotes the Picard stack of X over B
% \Picardfunctor_{X/B} denotes the Picard functor of X over B
\def\Pic{\mathop{\mathrm{Pic}}\nolimits}
\def\Picardstack{\mathcal{P}\!\mathit{ic}}
\def\Picardfunctor{\mathrm{Pic}}
\def\Deformationcategory{\mathcal{D}\!\mathit{ef}}

\setlength{\emergencystretch}{3em}
\begin{document}
\title[Stacks Project, Tag 0ECD]{327 --- Affineness of the maximal etale locus over excellent regular schemes}
\maketitle
\noindent Copyright (C) 2005--2025 Johan de Jong. Stacks Project authors. Source: \href{https://stacks.math.columbia.edu/tag/0ECD}{Tag 0ECD}.

\noindent Editorial difficulty note (not author text). Difficulty H3: The proof must show the etale locus is affine even when its complement is not a divisor. It localizes over the regular target, uses purity and Hartshorne-Lichtenbaum vanishing to bound the structure-sheaf cohomological dimension of the locus, then applies the prior local affineness criterion to finish the dimension induction..

\noindent Editorial provenance note (not author text). History: excerpt from revision \texttt{a0444\allowbreak 6e57e\allowbreak c1fbc\allowbreak 252a8\allowbreak 71afc\allowbreak ec775\allowbreak 2fb28\allowbreak 07b14}, pione.tex, lines 6512--6572. Reference presentation was adapted; original-author context and core support blocks were retained or added. The mathematical wording is unchanged. Licensed under GNU FDL 1.2 or later; no invariant sections or cover texts. See ../licenses/COPYING. Context blocks reproduce author definitions and supporting results; they are not additional counted problems.

\begin{theorem}
\label{theorem-global}
Let $Y$ be an excellent regular scheme over a field. Let $f : X \to Y$
be a finite type morphism of schemes with $X$ normal. Let $V \subset X$
be the maximal open subscheme where $f$ is \'etale. Then the inclusion
morphism $V \to X$ is affine.
\end{theorem}

\begin{proof}
Let $x \in X$ with image $y \in Y$. It suffices to prove that
$V \cap W$ is affine for some affine open neighbourhood $W$ of $x$.
Since $\Spec(\mathcal{O}_{X, x})$ is the limit of the schemes $W$,
this holds if and only if
$$
V_x = V \times_X \Spec(\mathcal{O}_{X, x})
$$
is affine (Limits, Lemma \href{https://stacks.math.columbia.edu/tag/01Z6}{lemma limit affine (Tag 01Z6)}).
Thus, if the theorem holds for the morphism
$X \times_Y \Spec(\mathcal{O}_{Y, y}) \to \Spec(\mathcal{O}_{Y, y})$,
then the theorem holds. In particular, we may assume $Y$
is regular of finite dimension, which allows us to do induction
on the dimension $d = \dim(Y)$. Combining this with the same argument again,
we may assume that $Y$ is local with closed point $y$ and that
$V \cap (X \setminus f^{-1}(\{y\}) \to X \setminus f^{-1}(\{y\})$
is affine.

\medskip\noindent
Let $x \in X$ be a point lying over $y$. If $x \in V$, then
there is nothing to prove. Observe that  $f^{-1}(\{y\}) \cap V$
is a finite set of closed points (the fibres of an \'etale morphism
are discrete). Thus after replacing $X$
by an affine open neighbourhood of $x$ we may assume
$y \not \in f(V)$. We have to prove that $V$ is affine.

\medskip\noindent
Let $e(V)$ be the maximum $i$ with $H^i(V, \mathcal{O}_V) \not = 0$.
As $X$ is affine the integer $e(V)$ is the maximum of the numbers $e(V_x)$
where $x \in X \setminus V$, see
Local Cohomology, Lemma \href{https://stacks.math.columbia.edu/tag/0DXB}{lemma cd local (Tag 0DXB)}
and the characterization of cohomological dimension
in Local Cohomology, Lemma \href{https://stacks.math.columbia.edu/tag/0DX7}{lemma cd (Tag 0DX7)}.
We have $e(V_x) \leq \dim(\mathcal{O}_{X, x}) - 1$ by
Local Cohomology, Lemma \href{https://stacks.math.columbia.edu/tag/0DXC}{lemma cd dimension (Tag 0DXC)}.
If $\dim(\mathcal{O}_{X, x}) \geq 2$ then purity of
ramification locus (Lemma \href{https://stacks.math.columbia.edu/tag/0EA4}{lemma purity ramification (Tag 0EA4)})
shows that $V_x$ is strictly smaller than the punctured spectrum of
$\mathcal{O}_{X, x}$. Since $\mathcal{O}_{X, x}$ is
normal and excellent, this implies
$e(V_x) \leq \dim(\mathcal{O}_{X, x}) - 2$ by
Hartshorne-Lichtenbaum vanishing
(Local Cohomology, Lemma \href{https://stacks.math.columbia.edu/tag/0EB7}{lemma affine complement (Tag 0EB7)}).
On the other hand, since $X \to Y$ is of finite type
and $V \subset X$ is dense (after possibly replacing $X$
by the closure of $V$), we see that $\dim(\mathcal{O}_{X, x}) \leq d$
by the dimension formula
(Morphisms, Lemma \href{https://stacks.math.columbia.edu/tag/02JU}{lemma dimension formula (Tag 02JU)}).
Whence $e(V) \leq \max(0, d -  2)$.
Thus $V$ is affine by Lemma \href{https://stacks.math.columbia.edu/tag/0ECC}{lemma conclude (Tag 0ECC)}
if $d \geq 2$. If $d = 1$ or $d = 0$, then the punctured spectrum of
$\mathcal{O}_{Y, y}$ is affine and hence $V$ is affine.
\end{proof}
\end{document}
